\documentclass[11pt,a4paper]{article}

\usepackage{amsmath,amssymb,amsthm}
\usepackage{mathtools}
\usepackage{hyperref}
\usepackage{booktabs}
\usepackage{enumitem}
\usepackage[margin=1.15in]{geometry}
\usepackage{xcolor}

% Shortcuts
\newcommand{\Gstar}{{G^{*}}}
\newcommand{\vpi}{\varpi}
\newcommand{\Nc}{N_c}
\newcommand{\Neff}{N_{\mathrm{eff}}}
\newcommand{\KB}{K_B}
\newcommand{\ZZ}{\mathbb{Z}}
\newcommand{\RR}{\mathbb{R}}
\newcommand{\QQ}{\mathbb{Q}}
\newcommand{\UU}{\mathcal{U}}
\newcommand{\OO}{\Omega}
\newcommand{\PP}{\Pi}
\newcommand{\onticell}{\ell}

% Inline epistemic tags
\newcommand{\etag}[1]{{\normalfont\textsc{\footnotesize [#1]}}}

% Theorem environments
\newtheoremstyle{compact}%
  {6pt}{6pt}{\itshape}{}{\bfseries}{.}{.5em}{}
\theoremstyle{compact}
\newtheorem{definition}{Definition}
\newtheorem{axiom}{Axiom}
\newtheorem{proposition}{Proposition}
\newtheorem{theorem}{Theorem}
\newtheorem{conjecture}{Conjecture}
\newtheorem*{remark}{Remark}

% Running header
\usepackage{fancyhdr}
\pagestyle{fancy}
\fancyhf{}
\fancyhead[R]{\small\itshape The Instantiating Formula}
\fancyfoot[C]{\thepage}
\renewcommand{\headrulewidth}{0.4pt}

\title{\textsc{The Instantiating Formula}\\[6pt]
\large An Ontological Constructor for\\
Foundational Ternary Dynamics}

\author{}
\date{}

\begin{document}
\maketitle
\thispagestyle{empty}

\begin{abstract}
\noindent We formalize the ontological update law of Foundational Ternary Dynamics (FTD) as a single constructor function. The \emph{instantiating formula} takes as input the cumulative record of all prior actualizations, evaluates the space of admissible continuations weighted by an action-coherence kernel, and projects the result through an integer selector to produce the next actualized state. The formula is defined using only the bridge constant $\Gstar = \sqrt{2}\,\Gamma(1/4)^2/(2\pi)$, the ontic ratio $\onticell = \vpi^2/{G^*}^2 = \pi/4$, and the framework integers $I = \{3, 4, 7, 13\}$. We show that in the lattice limit it reduces to the engine's tick cycle, in the static limit it reduces to energy minimization, and in the measurement limit it recovers the Gauss-constraint projection. The formula itself is classified as a \emph{conjecture}: it unifies the component results proved in the companion paper~\cite{lifecycle} but does not yet derive them. Falsification criteria are stated.
\end{abstract}

\smallskip

\begin{quote}\itshape
Every theorem is a tautology that took work to see. Every constructor is a tautology that takes a tick to execute.
\end{quote}


%=======================================================================
\section{Why a Constructor?}
\label{sec:why}
%=======================================================================

The companion paper~\cite{lifecycle} traces the lifecycle of a hydrogen atom through a chain of individually proved results: the master quadratic produces $\alpha$, the Moore decomposition produces the gauge groups, the Gauss constraint produces Bell violation, the lossy projection produces the arrow of time. Each result has its own proof and its own epistemic tag.

But the proofs are separate. They share the same constants---$\Gstar$, $\onticell$, $\{3, 4, 7, 13\}$---but no single law generates all of them. The question this paper asks is: \emph{is there a single update rule from which every component result follows as a special case?}

Standard physics faces the same question at a different level. The Standard Model Lagrangian unifies the electromagnetic, weak, and strong interactions into one expression. General relativity unifies geometry and gravity into one equation. Neither reduces to the other. FTD, if it is to be a framework rather than a collection of coincidences, requires a comparable unification---not of forces, but of the update process itself.

We propose such a law. It is classified as a \emph{conjecture} because it has not been derived from Axiom Zero. What we can show is that it is \emph{consistent} with every proved result: it reduces to them in appropriate limits. Whether it is the \emph{unique} law with this property is an open question.


%=======================================================================
\section{The Formal Framework}
\label{sec:framework}
%=======================================================================

We define each symbol in the constructor before assembling them.

\begin{definition}[Cumulative Record]
\label{def:record}
Let $\mathbf{L} = \ZZ^3$ be the cubic lattice. At each lattice site $\mathbf{x} \in \mathbf{L}$, define the \emph{ontic state} as the pair $(s(\mathbf{x}),\, \mathbf{J}(\mathbf{x}))$ where $s \in \{-1, 0, +1\}$ is the ternary state and $\mathbf{J} \in \RR^3$ is the flux vector. The \emph{instantiated universe at update index $k$} is
\[
\UU_k \;=\; \bigl\{(s_k(\mathbf{x}),\, \mathbf{J}_k(\mathbf{x}))\bigr\}_{\mathbf{x} \in \mathbf{L}}.
\]
The \emph{cumulative record} is the ordered sequence $\UU_{[0:k]} = (\UU_0,\, \UU_1,\, \ldots,\, \UU_k)$. \etag{Definition}
\end{definition}

\begin{remark}
The cumulative record is not a convenience. It is the ontology. The past is not erased; it is the substrate from which the future is projected. This is why the projection $\Pi$ is lossy (information is discarded in the $\UU_k$ snapshot but exists in $\UU_{[0:k]}$) and why time has a direction.
\end{remark}

\begin{definition}[Admissible Continuations]
\label{def:continuations}
Given $\UU_k$, the space of \emph{admissible continuations} is
\[
\OO(\UU_k) \;=\; \bigl\{\UU' : \UU' \text{ is reachable from } \UU_k \text{ by one Moore-neighborhood update}\bigr\}.
\]
Admissibility requires: (i)~locality---each site's new state depends only on its 26 neighbors; (ii)~ternary closure---$s' \in \{-1, 0, +1\}$; (iii)~CFL stability---no information propagates faster than $c = 1/\sqrt{3}$ lattice units per tick. \etag{Definition}
\end{definition}

\begin{definition}[Ontic Action]
\label{def:action}
The \emph{ontic action} of a continuation $\UU' \in \OO(\UU_k)$ is \etag{Theorem (from engine energy functional)}
\[
\mathcal{A}[\UU' \,|\, \UU_k] \;=\; \sum_{\mathbf{x} \in \mathbf{L}} \left[\frac{1}{2}\bigl|\mathbf{J}'(\mathbf{x}) - \mathbf{J}_k(\mathbf{x})\bigr|^2 \;+\; \KB\,\bigl|s'(\mathbf{x}) - s_k(\mathbf{x})\bigr|\right]
\]
where $\KB = 0.5100$\,MeV is the manifestation threshold (FTD-derived: $m_e = m_P\sqrt{2\pi}\,(16/3)\,\alpha^{11}$). The first term penalizes flux change (kinetic cost). The second penalizes state change (manifestation cost). A continuation that changes nothing has $\mathcal{A} = 0$.
\end{definition}

This is not a new object. It is the energy functional $\mathcal{E} = \frac{1}{2}|\mathbf{J}|^2 + \KB|s|$ from the engine specification~\cite{lifecycle}, rewritten as a cost of transition rather than a cost of existence. The two forms are related by
\[
\mathcal{A}[\UU' \,|\, \UU_k] \;=\; \mathcal{E}[\UU'] - \mathcal{E}[\UU_k] + \text{(boundary terms)}.
\]

\begin{remark}
The action $\mathcal{A}$ encodes dissipation implicitly. In the engine~\cite{engine}, Rayleigh damping at rate $\gamma = \alpha$ reduces flux amplitude by $(1-\alpha)$ per tick. The coupling amplitude is $g_c = \sqrt{\alpha}$, so $\gamma = g_c^2$---the dissipation rate equals the coupling amplitude squared, mirroring the QED relationship between transition rates and coupling constants. This identification is \etag{Imposed} in the engine; deriving it from Axiom Zero remains open. In the constructor, dissipation appears as the weighting $e^{-\onticell\,\mathcal{A}}$: continuations with large action are exponentially suppressed, and the suppression scale $\onticell = \vpi^2/{G^*}^2$ is fixed by the lemniscatic constants.
\end{remark}

\begin{definition}[Coherence Functional]
\label{def:coherence}
The \emph{coherence} of a continuation $\UU'$ measures how well it satisfies the structural constraints of the framework: \etag{Theorem (Gauss term) + Selection (Wilson term)}
\[
\mathcal{C}[\UU' \,|\, \UU_k] \;=\; \underbrace{\sum_{\mathbf{x}} \bigl(1 - |\nabla \cdot \mathbf{J}'(\mathbf{x}) - \rho'(\mathbf{x})|^2\bigr)}_{\text{Gauss constraint satisfaction}} \;+\; \underbrace{\sum_{R,T \geq 1} w_{RT}\,\ln\bigl\langle W(R,T)\bigr\rangle'}_{\text{Wilson area-law enforcement}}
\]
where $\rho'(\mathbf{x}) = s'(\mathbf{x})$ is the charge density and $\langle W(R,T)\rangle'$ is the expectation value of the $R \times T$ rectangular Wilson loop in the continuation $\UU'$.

\smallskip
\noindent\textit{The Gauss term.} Maximized when $\nabla \cdot \mathbf{J}' = \rho'$ everywhere. This eliminates one flux degree of freedom per site, projecting $\RR^3 \to \RR^2$ (the transverse plane), which is the mechanism producing $S = 2\sqrt{2}$ in the Bell analysis~\cite{lifecycle}.

\smallskip
\noindent\textit{The Wilson term.} On a compact U(1) lattice gauge theory at coupling $\beta$, the single-plaquette integral gives the exact area law~\cite{wilson_lgt}:
\[
\langle W(R,T)\rangle \;=\; u_p^{\,RT}, \qquad u_p \;=\; \frac{I_1(\beta)}{I_0(\beta)}
\]
where $I_n$ are modified Bessel functions of the first kind. At $\beta = x_- = 3.024$:
\[
u_p(x_-) = 0.812, \qquad \sigma(x_-) = -\ln u_p = 0.209.
\]
The string tension $\sigma = 0.209$ is an FTD output, not an input---it follows from the smaller root of the master quadratic. The weights $w_{RT} > 0$ are monotonically decreasing in $R + T$ (larger loops contribute less); the simplest choice is $w_{RT} = (RT)^{-1}$, which makes the sum converge and penalizes area-law violation at all scales.
\end{definition}

\begin{remark}
The Gauss term is rigorous: it is a quadratic penalty computable site-by-site. The Wilson term is a \emph{selection}: the area-law form is exact in the strong-coupling expansion of compact U(1), and the choice of weights $w_{RT} = (RT)^{-1}$ is natural but not uniquely forced. An alternative weight scheme would change the relative importance of short- vs.\ long-distance confinement but would not change the qualitative structure (area law at $x_-$, Coulomb law at $x_+$).
\end{remark}

\begin{definition}[The Selector-Projector]
\label{def:selector}
Let $I = \{3, 4, 7, 13\}$ be the framework integers. The \emph{selector-projector} $\PP_I$ is the map \etag{Selection}
\[
\PP_I: \mathcal{W}(\OO(\UU_k)) \;\longrightarrow\; \OO(\UU_k)
\]
that takes a weighted distribution over admissible continuations and returns a single actualized state by enforcing three successive constraints:

\smallskip
\noindent\textit{Step 1: Integer reduction.} Only $\Nc = 3$ is free. The remaining integers follow from independent physical arguments~\cite{lifecycle}:
\[
N_{\mathrm{base}} = |\mathrm{Aut}(E)| = 4, \qquad b_3 = \frac{11\Nc - 2N_f}{3} = 7 \;\;(N_f = 6), \qquad \Neff = b_3 + 2\Nc = 13.
\]
$N_{\mathrm{base}}$ is the automorphism order of the elliptic curve $E: y^2 = x^3 - x$ (SP1). $b_3$ is the one-loop QCD beta function coefficient for $\Nc = 3$ colors and $N_f = 6$ quark flavors. $\Neff$ is their derived combination. Once $\Nc$ is given, these are determined.

\smallskip
\noindent\textit{Step 2: Gauss projection.} The Gauss constraint $\nabla \cdot \mathbf{J} = \rho$ eliminates the longitudinal flux component at each site, reducing $\RR^3 \to \RR^2$. This is a linear projection and therefore unique: the transverse subspace $\mathbf{J}_{\mathrm{phys}} = \mathbf{J} - \hat{\mathbf{k}}(\hat{\mathbf{k}} \cdot \mathbf{J})$ is uniquely determined.

\smallskip
\noindent\textit{Step 3: Ternary crystallization.} The continuous flux is mapped to discrete ternary states via the sign function: $s(\mathbf{x}) = \mathrm{sgn}(\mathbf{J}_{\mathrm{phys}}(\mathbf{x}) \cdot \hat{\mathbf{n}})$ for a locally determined axis $\hat{\mathbf{n}}$. This is an irreversible coarse-graining that maps $\RR^2 \to \{-1, 0, +1\}$.

\smallskip
\noindent\textit{Uniqueness.} Each step is deterministic: Step~1 is algebraic, Step~2 is a linear projection (unique kernel), Step~3 is a pointwise function. The composition $\PP_I = \mathrm{sgn} \circ \Pi_{\mathrm{Gauss}} \circ \mathrm{reduce}_I$ is therefore deterministic. Given the weighted distribution $\mathcal{Z}$, the output is unique.

The projection is lossy: the codomain $\{-1,0,+1\}^{|\mathbf{L}|}$ has cardinality $3^{(N^3)}$ while the domain is infinite-dimensional. This information loss is the source of the arrow of time.
\end{definition}

\begin{definition}[Ontic Adjunction]
\label{def:adjunction}
The operation $\oplus$ appends an actualized state to the cumulative record: \etag{Definition}
\[
\UU_{[0:k+1]} \;=\; \UU_{[0:k]} \oplus \UU_{k+1}
\]
This is concatenation, not superposition. The record grows monotonically.
\end{definition}


%=======================================================================
\section{The Formula}
\label{sec:formula}
%=======================================================================

\begin{conjecture}[The Instantiating Formula]
\label{conj:formula}
Let $\onticell = \vpi^2/{G^*}^2 = \pi/4$ be the ontic ratio. The next instantiated state is given by the constructor
\begin{equation}
\label{eq:constructor}
\boxed{
\UU_{k+1} \;=\; \PP_I\!\left[\,\UU_k \;;\; \mathcal{Z}(\cdot\,|\,\UU_k)\,\right]
}
\end{equation}
where the \emph{instantiating kernel} is the weight function over admissible continuations:
\begin{equation}
\label{eq:kernel}
\mathcal{Z}(\UU'\,|\,\UU_k) \;=\; \exp\!\left(-\onticell\,\mathcal{A}[\UU'\,|\,\UU_k] \;+\; \frac{1}{{G^*}^2}\,\mathcal{C}[\UU'\,|\,\UU_k]\right).
\end{equation}
The selector-projector $\PP_I$ receives the full weighted distribution $\{\mathcal{Z}(\UU'\,|\,\UU_k)\}_{\UU' \in \OO(\UU_k)}$ and returns the unique actualization consistent with the integer constraints $I = \{3, 4, 7, 13\}$.
\end{conjecture}

\subsection*{Reading the Formula}

The formula says: at each tick, reality evaluates every admissible continuation of itself, weights each by how much it costs (action) and how well it coheres (coherence), and projects the weighted distribution through the integer selector to produce one actualized outcome.

The three terms do different work:
\begin{itemize}[nosep]
\item $\onticell\,\mathcal{A}$: penalizes change. High action $\to$ low weight. This is the principle of least action, discretized.
\item ${G^*}^{-2}\,\mathcal{C}$: rewards coherence. High coherence $\to$ high weight. This is the Gauss constraint and topological conservation, enforced softly.
\item $\PP_I$: quantizes the outcome. Continuous weights become a discrete actuality. This is where the bridge constant earns its name.
\end{itemize}

\subsection*{Connection to Statistical Mechanics}

The kernel $\mathcal{Z} = \exp(-\onticell\,\mathcal{A} + {G^*}^{-2}\,\mathcal{C})$ is a Boltzmann weight. In statistical mechanics notation, $\exp(-\beta E + \beta\mu N)$ weights configurations by energy $E$ at inverse temperature $\beta$ and chemical potential $\mu$. The constructor maps onto this structure:

\begin{center}
\begin{tabular}{lcl}
\toprule
\textbf{Statistical mechanics} & & \textbf{Constructor} \\
\midrule
Inverse temperature $\beta$ & $\longleftrightarrow$ & Ontic ratio $\onticell = \pi/4$ \\
Energy $E$ & $\longleftrightarrow$ & Ontic action $\mathcal{A}$ \\
Chemical potential $\beta\mu$ & $\longleftrightarrow$ & ${G^*}^{-2}$ \\
Conserved charge $N$ & $\longleftrightarrow$ & Coherence $\mathcal{C}$ \\
Partition function $Z$ & $\longleftrightarrow$ & $\sum_{\UU'} \mathcal{Z}(\UU'\,|\,\UU_k)$ \\
\bottomrule
\end{tabular}
\end{center}

This is not a metaphor. The lattice partition function $Z = \sum \exp(-S_E)$ with Euclidean action $S_E$ is the standard starting point for lattice QFT. The constructor's kernel is this partition function, with the action decomposed into a cost term $\mathcal{A}$ and a constraint-satisfaction term $\mathcal{C}$.

The identification has a concrete consequence: the ``temperature'' of reality is not a free parameter. It is $\onticell^{-1} = 4/\pi = {G^*}^2/\vpi^2$---a fixed ratio of the lemniscatic constants. There is no thermal equilibrium to reach; the constructor operates at a single, determined temperature. This is consistent with the deterministic axiom: a system at fixed temperature with fixed constraints has a unique Gibbs state.

\subsection*{Why Not Extremization?}

One might expect the constructor to select the single highest-weighted continuation---the $\operatorname{arg\,max}$ of $\mathcal{Z}$. This would be incorrect for two reasons.

First, a weighted sum over configurations does not select a single maximum. It integrates over \emph{all} configurations, each contributing proportionally to its weight. The dominant configuration prevails in the classical limit but not in general. The correct formulation is the full weighted distribution, which $\PP_I$ receives.

Second, the projection $\PP_I$ is not maximization. It is \emph{quantization}: the continuous distribution is mapped to a discrete outcome by the integer constraints, just as the Gauss constraint maps three-dimensional flux to its two-dimensional transverse component. The information lost in this projection is the arrow of time.


%=======================================================================
\section{Recovery of Known Results}
\label{sec:recovery}
%=======================================================================

If the constructor is correct, every proved result from~\cite{lifecycle} must appear as a special case.

\begin{proposition}[Static Limit]
\label{prop:static}
If $\UU_k$ is a ground state ($\mathcal{A} = 0$ for the identity continuation), the constructor reduces to
\[
\UU_{k+1} \;=\; \PP_I\!\left[\UU_k \;;\; e^{{G^*}^{-2}\,\mathcal{C}_{\max}}\right] \;=\; \UU_k.
\]
The ground state is a fixed point of the constructor. \etag{Theorem}
\end{proposition}

\begin{proof}
If $\UU' = \UU_k$ is admissible and has $\mathcal{A} = 0$ and maximal $\mathcal{C}$, it dominates the kernel. The projection returns $\UU_k$ itself. Stationarity is energy minimization.
\end{proof}

\begin{conjecture}[Lattice Limit]
\label{prop:lattice}
On a finite lattice $\mathbf{L}_N = \{0, \ldots, N-1\}^3$ with the Moore neighborhood, the constructor reduces to the engine tick cycle:
\begin{align*}
&\text{phase\_read:} && \text{evaluate } \mathcal{Z}(\UU'\,|\,\UU_k) \text{ for all } \UU' \in \OO(\UU_k) \\
&\text{phase\_write:} && \text{select } \UU_{k+1} = \PP_I[\,\cdot\;] \\
&\text{gauss\_project:} && \text{enforce } \nabla \cdot \mathbf{J} = \rho \text{ (coherence maximization)} \\
&\text{phase\_forces:} && \text{compute flux gradients (action evaluation)} \\
&\text{phase\_movement:} && \text{propagate states (adjunction } \oplus\text{)}
\end{align*}
\end{conjecture}

This correspondence is not proved---it is the \emph{content} of the conjecture. The claim is that the five engine phases are a computational decomposition of the single formula~\eqref{eq:constructor}. Verifying this requires showing that the engine's output is independent of decomposition order, which is guaranteed by the deterministic update axiom but has not been formally connected to the kernel formulation.

\begin{proposition}[Measurement Limit]
\label{prop:measurement}
When the coherence functional is dominated by its Gauss-constraint term and the action is negligible ($\onticell\,\mathcal{A} \ll {G^*}^{-2}\,\mathcal{C}$), the constructor reduces to the Gauss projection: \etag{Theorem}
\[
\PP_I\!\left[\,\cdot\;;\; e^{{G^*}^{-2}\,\mathcal{C}}\,\right] \;\longrightarrow\; \mathbf{J}_{\mathrm{phys}} = \mathbf{J} - \hat{\mathbf{k}}(\hat{\mathbf{k}} \cdot \mathbf{J})
\]
which produces the Gauss projection $\RR^3 \to \RR^2$ (the transverse flux plane). A local sign-function measurement model on this 2D plane gives the triangle-wave correlation $E(\theta) = 1 - 2|\theta|/\pi$ and $S = 2$ (Bell bound). The engine's $S = 2\sqrt{2}$ result arises from the \emph{global} Poisson solver enforcing the constraint non-locally---see~\cite{lifecycle}, \S9 for the full analysis.
\end{proposition}

\begin{proof}
When $\mathcal{C}$ dominates, the kernel selects continuations with maximal Gauss-constraint satisfaction. The Gauss constraint $\nabla \cdot \mathbf{J} = \rho$ eliminates one flux degree of freedom, projecting $\RR^3 \to \RR^2$. The transverse projection is a theorem. The Bell violation $S = 2\sqrt{2}$ is a computational result of the non-local constraint solver~\cite{lifecycle}.
\end{proof}

\medskip
\noindent\textbf{Open Problem: Quadratic Recovery.}
\label{prop:quadratic}
The ontic ratio $\onticell$ and bridge constant $\Gstar$ appear in the kernel~\eqref{eq:kernel} as the two natural scales. Their balance defines a critical surface:
\[
\onticell\,\mathcal{A}[\UU'\,|\,\UU_k] \;=\; {G^*}^{-2}\,\mathcal{C}[\UU'\,|\,\UU_k] \qquad \text{(action cost equals coherence reward)}.
\]
The master quadratic $x^2 - 16{G^*}^2 x + 16{G^*}^3 = 0$ is proved in~\cite{lifecycle} from the elliptic curve $E: y^2 = x^3 - x$ and its arithmetic invariants. Its roots $x_+ = \alpha^{-1}$ and $\lfloor x_- \rfloor = \Nc = 3$ set the electromagnetic and strong coupling scales.

\emph{The open question is whether the constructor requires this specific quadratic, or merely accommodates it.} If the balance condition on the critical surface can be shown to produce the master quadratic---if the only polynomial consistent with the kernel's symmetries and the Watson integral is the one already proved---then the constructor would derive the coupling constants rather than inheriting them. This is not proved. It is the central open problem of this paper, and we do not disguise it as a result. \etag{Open}


%=======================================================================
\section{Consequences}
\label{sec:consequences}
%=======================================================================

If the constructor is accepted as a conjecture, several interpretive consequences follow.

\subsection*{1. Time Is Not Fundamental}

The update index $k$ is not a coordinate. It is the \emph{cardinality of the cumulative record}. Time is defined as:
\[
t \;\stackrel{\text{def}}{=}\; |\UU_{[0:k]}| \;=\; k + 1.
\]
There is no background clock. The constructor \emph{is} the clock. Each application of~\eqref{eq:constructor} increments $k$ by one. The ``flow of time'' is the iteration of the constructor.

\subsection*{2. Discrete Objects Are Selector-Stabilized Excitations}

A particle---a proton, an electron, a hydrogen atom---is a pattern in $\UU_k$ that the selector-projector $\PP_I$ returns to itself under repeated application. It is a \emph{fixed point} of $\PP_I$, not of the full constructor (which advances the clock) but of the structural content:
\[
\PP_I[\UU_k] \cong \PP_I[\UU_{k+1}] \qquad \text{(structural isomorphism, not identity)}.
\]
Stability means the integer constraints $I = \{3,4,7,13\}$ lock the pattern in place. Instability means the pattern eventually fails to satisfy the constraints and dissolves---as when a hydrogen atom is ionized or a proton decays.

\subsection*{3. Continuity Is Ontically Prior}

The space of admissible continuations $\OO(\UU_k)$ is continuous (parametrized by the flux field $\mathbf{J} \in \RR^3$ at each site). The actualized state $\UU_{k+1}$ is discrete (ternary states $s \in \{-1, 0, +1\}$ selected by $\PP_I$). The continuous is not an approximation to the discrete. The continuous is the \emph{possibility space} from which the discrete is projected.

This places FTD in an unusual ontological position. In standard quantum mechanics, the continuum (Hilbert space) is fundamental and discreteness (measurement outcomes) is imposed by observation. In Bohmian mechanics, a continuous guiding wave pilots a discrete particle. In digital physics (Wolfram, Fredkin), everything is discrete and the continuum is emergent.

FTD is none of these. It is \emph{ontologically discrete} (the lattice is the substrate) but \emph{explanatorily continuous} (the flux field is dispositional---it exists ``before'' the ternary state in the order of explanation). The constructor makes this hierarchy explicit: $\OO(\UU_k)$ is continuous, $\mathcal{Z}$ is real-valued, $\PP_I$ is discrete. The continuous layer is not emergent. It is the necessary precondition for the discrete layer to be selected. Without the flux, there is nothing to project.

The closest parallel in existing philosophy is \emph{ontic structural realism}: what is fundamental is neither things nor properties but \emph{structure}. In FTD, the structure is the constructor itself---the three-layer map from possibility through weighting to actualization. Particles are not fundamental objects; they are stable patterns in the constructor's output. Fields are not fundamental objects; they are the constructor's input space. What is fundamental is the map.

\subsection*{4. $\Gstar$ Performs the Bookkeeping}

The bridge constant $\Gstar$ appears in the kernel in two places:
\begin{itemize}[nosep]
\item In $\onticell = \vpi^2/{G^*}^2$: weighting the action cost.
\item In ${G^*}^{-2}$: weighting the coherence reward.
\end{itemize}
Both are functions of $\Gstar$ alone. The constructor has no free parameters beyond the framework integers and the bridge constant. Everything---the strength of action penalization, the reward for coherence, the scales at which coupling constants emerge---is fixed by one transcendental number and four integers.

\subsection*{5. The Ontic Ratio $\onticell$ Is the Primitive}

Since $\onticell = \vpi^2/{G^*}^2 = \pi/4$, the number $\pi$ is a derived quantity:
\[
\pi = 4\onticell.
\]
Within the constructor, $\pi$ does not appear independently. It appears only through $\onticell$, which governs the tension between action (cost of change) and the circular geometry implicit in the lemniscatic constants. The hierarchy is: $\vpi$ and $\Gstar$ are primitive; $\onticell$ is their ratio; $\pi$ is four times $\onticell$.

\subsection*{6. The Absolute Mass Scale Is Selected}

The constructor has no free parameters. But does it determine the absolute mass scale?

In lattice natural units ($\KB = 1$, $a = 1$, $\tau = 1$), every mass is a dimensionless number:
\begin{center}
\begin{tabular}{lrl}
\toprule
\textbf{Particle} & \textbf{Mass (lattice units)} & \textbf{Formula} \\
\midrule
Electron & $1$ & $\KB$ (by definition) \\
Proton & $1836.47$ & $\Neff/\alpha + T(10)$ \\
Higgs & $2.46 \times 10^5$ & $v\sqrt{6/23}\,/\,\KB$ \\
Planck & $2.39 \times 10^{22}$ & $1/(\sqrt{2\pi}\cdot 16/3 \cdot \alpha^{11})$ \\
\bottomrule
\end{tabular}
\end{center}
All entries are computable from $\Gstar$ and the framework integers. The Planck mass is not an input---it is the ratio of the manifestation scale to the gravitational scale, both determined by the constructor's constants.

The identification $\KB = m_e$ follows from a four-step chain: the Born--Infeld Lagrangian assigns rest energy $\KB$ to each manifested voxel \etag{Axiom}; manifestation creates charge \etag{Theorem}; one voxel is the minimum charged excitation \etag{Theorem}; the minimum charged excitation is the electron \etag{Selection}. The ``one external calibration'' is the unit convention $\KB = 0.5100$~MeV---a choice of measurement units, not a physical parameter.


%=======================================================================
\section{Compressed Form}
\label{sec:compressed}
%=======================================================================

The constructor may be written in one line:
\begin{equation}
\label{eq:compressed}
\boxed{
\UU_{k+1} \;=\; \PP_I\!\left[\,\UU_k \;;\; e^{-\onticell\,\mathcal{A} \,+\, {G^*}^{-2}\,\mathcal{C}}\,\right]
}
\end{equation}
or in words:
\begin{center}
\framebox{\parbox{0.85\textwidth}{\centering
\textit{The next state of reality is the integer-projected outcome\\
of the action-coherence weighted evaluation of all admissible continuations\\
of the current cumulative record.}
}}
\end{center}

\medskip

The constructor has three layers:
\begin{center}
\begin{tabular}{lcl}
\toprule
\textbf{Layer} & \textbf{Symbol} & \textbf{Ontological role} \\
\midrule
Possibility & $\OO(\UU_k)$ & What could happen next (continuous) \\
Weighting & $\mathcal{Z} = e^{-\onticell\mathcal{A} + {G^*}^{-2}\mathcal{C}}$ & How costly and how coherent (real-valued) \\
Actualization & $\PP_I$ & What does happen (discrete) \\
\bottomrule
\end{tabular}
\end{center}

This three-layer structure mirrors the two-layer ontology of Axiom Zero (flux/state $\equiv$ disposition/actuality) with the addition of an evaluation layer that mediates between them.


%=======================================================================
\section{Epistemic Status}
\label{sec:epistemic}
%=======================================================================

\subsection*{What Is Proved}

\begin{center}
\small
\begin{tabular}{lll}
\toprule
\textbf{Claim} & \textbf{Status} & \textbf{Where proved} \\
\midrule
Definitions 1--6 (formal objects) & \etag{Definition} & This paper, \S\ref{sec:framework} \\
Static limit (Prop.~\ref{prop:static}) & \etag{Theorem} & This paper, \S\ref{sec:recovery} \\
Measurement limit (Prop.~\ref{prop:measurement}) & \etag{Theorem} & This paper + \cite{lifecycle} \S9 \\
Gauss $\to$ Bell violation & \etag{Theorem} & \cite{lifecycle} \S9 \\
Projection is lossy $\to$ arrow of time & \etag{Theorem} & \cite{lifecycle} \S10 \\
\bottomrule
\end{tabular}
\end{center}

\subsection*{What Is Conjectured}

\begin{center}
\small
\begin{tabular}{lll}
\toprule
\textbf{Claim} & \textbf{Status} & \textbf{What would prove it} \\
\midrule
The constructor itself (Conj.~\ref{conj:formula}) & \etag{Conjecture} & Derivation from Axiom Zero \\
Lattice limit (Conj.~\ref{prop:lattice}) & \etag{Conjecture} & Show engine $=$ constructor on finite $\mathbf{L}$ \\
Quadratic recovery (\S\ref{sec:recovery}) & \etag{Open} & Derive master quadratic from balance condition \\
Coherence functional (Def.~\ref{def:coherence}) & \etag{Selection} & Prove uniqueness of weight scheme $w_{RT}$ \\
\bottomrule
\end{tabular}
\end{center}

\subsection*{What This Paper Does Not Claim}

\begin{enumerate}[nosep]
\item The constructor is \textbf{not derived} from Axiom Zero. It is a proposed unification of separately proved results.
\item The constructor does \textbf{not make new predictions} beyond those already made by the component results. It is an interpretive framework, not a computational advance.
\item The coherence functional $\mathcal{C}$ has \textbf{two tiers of certainty}. Its Gauss-constraint term is rigorous. Its Wilson area-law term uses exact Bessel-function formulas but depends on a weight scheme ($w_{RT} = (RT)^{-1}$) that is natural but not uniquely forced.
\item The formula does \textbf{not replace} the engine tick cycle as a computational tool. The engine is exact and implementable; the constructor is a conjecture about what the engine computes.
\end{enumerate}

\subsection*{Falsification}

The constructor inherits all falsification criteria from~\cite{lifecycle}. It adds one specific to its own structure:

\begin{enumerate}[nosep]
\item[\textbf{F7.}] \textbf{Non-equivalence with engine.} If a finite-lattice computation shows that the constructor~\eqref{eq:constructor} and the engine tick cycle produce different outputs for any initial state, the constructor is falsified as a description of the engine's dynamics.
\end{enumerate}

This is testable in principle by exhaustive enumeration on small lattices ($N \leq 8$).


%=======================================================================
\section*{Coda}
%=======================================================================

The instantiating formula is not a discovery. It is a conjecture about the form of the law that the component results obey. If the conjecture is correct, then reality is an iterated constructor: at each tick, the continuous possibility space is weighted by action and coherence, projected through integer constraints, and appended to the cumulative record.

The constructor has no free parameters. Its scales are set by one transcendental number ($\Gstar$) and four integers ($\{3, 4, 7, 13\}$). Its structure is three-layered: possibility (continuous), weighting (real-valued), actualization (discrete). Its iteration is time. Its lossy projection is entropy. Its fixed points are particles.

Whether this is the correct unification of the component results, or merely one consistent interpretation among others, is for future work to determine. The constructor is stated. Its consequences are explicit. Its assumptions are tagged. It is offered not as a conclusion but as a hypothesis---the hypothesis that existence has a formula, and the formula looks like this.

\medskip

\begin{center}\itshape
The bridge is not $\pi$. The bridge is $\Gstar$.\\[3pt]
The constructor is not the world.\\
The constructor is a conjecture about the world.
\end{center}


%=======================================================================
\begin{thebibliography}{99}
%=======================================================================

\bibitem{lifecycle} ``One Unit of Existence: The Lifecycle of a Hydrogen Atom from Lattice Geometry to Dissolution,'' companion paper (2026).

\bibitem{codata} E.~Tiesinga \textit{et al.}, ``CODATA recommended values of the fundamental physical constants: 2022,'' \textit{Rev.\ Mod.\ Phys.}\ \textbf{96} (2024), 015002.

\bibitem{deser} S.~Deser, ``Self-interaction and gauge invariance,'' \textit{Gen.\ Relativ.\ Gravit.}\ \textbf{1} (1970), 9--18.

\bibitem{watson} G.~N.~Watson, ``Three triple integrals,'' \textit{Q.\ J.\ Math.}\ \textbf{10} (1939), 266--276.

\bibitem{silverman} J.~H.~Silverman, \textit{The Arithmetic of Elliptic Curves}, 2nd ed., Springer (2009).

\bibitem{bell} J.~S.~Bell, ``On the Einstein--Podolsky--Rosen paradox,'' \textit{Physics}\ \textbf{1} (1964), 195--200.

\bibitem{tsirelson} B.~S.~Tsirelson, ``Quantum generalizations of Bell's inequality,'' \textit{Lett.\ Math.\ Phys.}\ \textbf{4} (1980), 93--100.

\bibitem{wilson_lgt} K.~G.~Wilson, ``Confinement of quarks,'' \textit{Phys.\ Rev.\ D}\ \textbf{10} (1974), 2445--2459.

\bibitem{chain} ``The Ontic Derivation Chain: Twenty-Six Physical Constants from Two Inputs,'' companion paper (2026).

\bibitem{bridge} ``The Discrete--Continuous Bridge Through the Lemniscatic Lens,'' companion paper (2026).

\bibitem{engine} ``The Lattice Engine: A Computational Implementation of Foundational Ternary Dynamics,'' companion paper (2026).

\end{thebibliography}

\end{document}
